% Current v17 frontier. Stable shape IDs; no existence assertion.
\subsection{残る十一型と，再開する具体的な入口}
\label{ns17:support-agenda}

現在の高分母 support の候補は\emph{十一型・三百五十六標識}である．
$G_{4,3}=K_4$ の二例外・四例外が解決し，その五標識を旧十二型から除いた．
再開時には最初に \nolinkurl{00_CONTINUATION_POLICY_ja.txt} を読む．
公開稿は，構成・幾何，五次整数接着，局所置換の三篇に分けた．
本節の graph 条件と局所算術の証明は五次整数接着論文で追い，
幾何的な atlas の実現と別種の core への置換は，それぞれ対応する稿で仮定を確認する．
以下の $G_{n,j}$ の識別子は維持し，欠番を詰めない．
全候補が五頂点全てを使い，$ij$ は無向辺 $\{i,j\}$ を表す．
$\mathbf d$ は降順の次数列，$\beta_1=|E|-4$ は cycle rank，
$L$ は葉の数，$N$ は $K_5$ 上の標識付き個数である．

\begin{center}
\footnotesize
\setlength{\tabcolsep}{3pt}
\begin{tabular}{crclrrr}
型 & $|E|$ & $\mathbf d$ & $E$ & $\beta_1$ & $L$ & $N$\\\hline
$G_{5,5}$ & 6 & $(33222)$ & $\{01,02,03,12,14,34\}$ & 2 & 0 & 60 \\
$G_{5,7}$ & 6 & $(33321)$ & $\{01,02,03,12,13,24\}$ & 2 & 1 & 60 \\
$G_{5,9}$ & 6 & $(43221)$ & $\{01,02,03,04,12,13\}$ & 2 & 1 & 60 \\
$G_{5,10}$ & 7 & $(33332)$ & $\{01,02,03,12,13,24,34\}$ & 3 & 0 & 30 \\
$G_{5,11}$ & 7 & $(43322)$ & $\{01,02,03,04,12,13,24\}$ & 3 & 0 & 60 \\
$G_{5,12}$ & 7 & $(43331)$ & $\{01,02,03,04,12,13,23\}$ & 3 & 1 & 20 \\
$G_{5,13}$ & 7 & $(44222)$ & $\{01,02,03,04,12,13,14\}$ & 3 & 0 & 10 \\
$G_{5,14}$ & 8 & $(43333)$ & $K_5-\{14,23\}$ & 4 & 0 & 15 \\
$G_{5,15}$ & 8 & $(44332)$ & $K_5-\{24,34\}$ & 4 & 0 & 30 \\
$G_{5,16}$ & 9 & $(44433)$ & $K_5-\{34\}$ & 5 & 0 & 10 \\
$G_{5,17}$ & 10 & $(44444)$ & $K_5$ & 6 & 0 & 1 \\
\end{tabular}
\end{center}

全 $1024$ 部分 graph と自己同型群を
\nolinkurl{verify_support_agenda_v17.py} で再検算している．
表は \nolinkurl{support_agenda_verification_v17.json} の \texttt{shapes} から再生成できる．
非空の support は全五頂点で連結であり，三角形を含み，$\beta_1\ge2$，辺数六以上である．
一つの core の四分母が全て一なら，support は空となる．
表は必要条件を通過した候補であり，整数接着や atlas の同型類の数ではない．

\paragraph{第一候補：$G_{5,5}$ と $G_{5,10}$ の共通低 puncture．}
$G_{5,5}$ の二価頂点 $3,4$ は，高分母近傍の間が低辺なので正の枝に限られる．
頂点 $2$ の近傍 $0,1$ は高辺で結ばれるため，その負の枝を先に捨てない．
高辺と共通低 puncture の組 $(ij;k)$ は
\[
 (02;4),\qquad(12;3),\qquad(34;2)
\]
であり，通常接着なら reverse の絶対値は四以下である．
最大約数に対応する例外接着を残して，三つの二高分母 core の低 row と node を接続する．
$G_{5,10}$ では二価頂点 $4$ が正の枝に限られ，
共通低 puncture の組は $(01;4)$ である．
一般の node 補題は中心の高分母次数や対向 $D$ の相殺を仮定しない．
一般三高分母分類を先に完成させる必要があると決めつけず，
固定した低 row と column からの二次式による局所復元を併用する．
代表 graph や標識を変える場合は，表示した組も再計算する．

\paragraph{第二候補：一高分母 hub への方法の拡張．}
$G_{5,12}$ は，$\{0,1,2,3\}$ 上の $K_4$ に高辺 $04$ を加えた型である．
第五 core $4$ の分母は $(q_{04},1,1,1)$ となる．
既知の一高分母 chart を使い，三本の低 hub 辺の node 条件と，
唯一の高 hub 辺の marked clutching・$D$ 相殺を先に組み合わせる．
$K_4$ 証明の全低 hub では全 pair product が零だったが，
この hub では高 ray と三つの低 rays の pair product は $q_{04}-1$ となる．
従って既存の零積条件をそのまま移してはいけない．
他の葉付き候補 $G_{5,7},G_{5,9}$ でも，一高分母 core とその隣接 core の接続条件が入口となる．
$K_4$ を部分 graph として含むだけで不可能と結論することはできない．

\paragraph{解決した $K_4$ から保存する手法．}
二例外では大きい負 reverse の端点を選び，正例外の三つの位置を分ければ，
hub の row/pair・node・$D$ 条件だけで矛盾する．
四例外では同一行の反対符号の大きい reverse の最小絶対値を $h$ とし，
$h=2,4$ と全領域 $h\ge6$ の三型に帰着した．
$h\ge6$ では hub entry の厳密符号が固定し，hub 条件が $h$ に依存しなくなる．
これは分母を切る探索ではなく，無限領域全体の同値な有限型表示である．
残る六候補には
\[
 (q_a,q_b,q_c)=(5,3,3),\quad
 (r_a,r_b,r_c)=(-4,1,2),\quad
 (x_a,x_b,x_c)=(-1,2,1)
\]
が共通し，局所 row/pair 条件は $6A^2+17A+14=0$ を要求する．
判別式 $-47$ が全候補を実数段階で排除する．
符号が安定する無限領域の有限化と，固定低 row/column からの二次式復元を他の型へ再利用する．
$K_4$ 自体の二例外・四例外は解決済みであり，
その未処理部分として循環変数を再探索する必要はない．

\paragraph{他の型と局所理論．}
三高分母の正符号 graph が常に matching となるか，
全ての高 $D$ が非負となり得るかは引き続き未解決である．
特異な rank 二 Schur 系を除外せず，既知の固定 slice の無限族も保持する．
$G_{5,14}=K_{2,2,1}$ では dense mod 五の二模型の整数 lifting が別の入口である．
高分母 support $H=\{ij:q_{ij}>1\}$ と素数 support $G_\ell=\{ij:\ell\mid q_{ij}\}$ を区別する．
局所整数族の存在から五 core の同時接着を結論しない．

二葉付き三角形の十模型・五十二組と $B$ 表現の七組も解決済みであり，再探索しない．
前者と後者は重複するので個数を加算しない．
共通低 puncture の node 上界，単位 entry の $(r,s,E)$ 復元，
有限な node 配分から中心三行を復元する手順を，残る型へ再利用する．
優先順位は固定せず，数学的な面白さ，インパクト，幾何的構成への寄与に応じて変更する．
